\documentclass[11pt,letterpaper]{article}
\usepackage[T1]{fontenc}
\usepackage[utf8]{inputenc}
\usepackage{lmodern}
\usepackage{amsmath,amssymb,amsthm}
\usepackage[margin=1in]{geometry}
\usepackage{microtype}
\usepackage{booktabs,longtable,array}
\usepackage{xcolor}
\usepackage{listings}
\usepackage{enumitem}
\usepackage[hidelinks]{hyperref}
\usepackage{fancyhdr}
\usepackage{lastpage}
\setlength{\headheight}{14pt}
\pagestyle{fancy}
\fancyhf{}
\fancyhead[L]{\small Uniform successor evaluation for iPRA}
\fancyhead[R]{\small Research report}
\fancyfoot[C]{\small \thepage}
\setlength{\parskip}{4pt}
\setlength{\parindent}{0pt}
\setlist{itemsep=2pt,topsep=5pt}
\emergencystretch=2em
\newtheorem{theorem}{Theorem}[section]
\newtheorem{lemma}[theorem]{Lemma}
\newtheorem{proposition}[theorem]{Proposition}
\newtheorem{corollary}[theorem]{Corollary}
\theoremstyle{definition}
\newtheorem{definition}[theorem]{Definition}
\newtheorem{remark}[theorem]{Remark}
\lstset{basicstyle=\ttfamily\scriptsize,columns=fullflexible,
  keepspaces=true,breaklines=true,showstringspaces=false,
  frame=single,rulecolor=\color{black!20},xleftmargin=4pt,xrightmargin=4pt}
\hypersetup{pdftitle={One successor step for equality and computation},
  pdfauthor={Research prepared for Vladimir Reshetnikov},
  pdfsubject={Five-rule clausal presentation of intuitionistic primitive recursive arithmetic}}
@@METRIC_MACROS@@
\begin{document}
\begin{titlepage}
\vspace*{1.0cm}
{\large Mathematical foundations \hfill 11 November 2026\par}
\vspace{1.2cm}
{\LARGE\bfseries One successor step for\\[4pt]
equality and computation\par}
\vspace{0.6cm}
{\Large A five-rule clausal presentation of\\
intuitionistic primitive recursive arithmetic\par}
\vspace{0.9cm}
{\large Prepared for Vladimir Reshetnikov\par}
\vspace{1.0cm}
\begin{abstract}
We give a simpler arithmetic inference interface for the current
\texttt{ConciseIPRA} proposal. On its already available zero-normal
iterator syntax, reflexivity and one uniform successor-evaluation rule
derive both computation and arbitrary equational transport. The three
syntactic alternatives of the current computational-equality rule are
replaced by one. The key is an explicit two-step derivation of equality
composition using a one-step identity iterator.

Together with the existing clausal logical interface, the complete
presentation has five primitive rule schemes: Use, Refl, Step, DE, and
DI. Its formulas are only equations and one two-variable binder.
We prove exact consequence equivalence with the native kernel and
therefore whole-domain bi-interpretability and equiconsistency with
full first-order iPRA. Both numerical composites have the identity
graph \(x=y\).

The fully validated standalone Python checker measures
\StepBytes\ normalized bytes, \StepTokens\ significant tokens, and
\StepAST\ AST nodes. A matched zero-normal computation/transport
control measures \ControlBytes, \ControlTokens, and \ControlAST,
respectively. The comparison counts normalization, syntax validation,
scope handling, clause storage, and the command-line interface.
The contribution is a local mathematical and implementation reduction;
no absolute minimum is asserted. The new equivalence is a written
proof supported by executable proof translations and boundary tests,
not a new Lean or Rocq verification.
\end{abstract}
\vfill
\textbf{Frozen comparison.}\quad
\texttt{VladimirReshetnikov/ProveIt}, commit\\
\texttt{6ca328039cafa1537d81ee6b67ee2dc54ffa5a4e}.\\
\textbf{Scope.}\quad Full intuitionistic first-order arithmetic,
including quantified formulas and arbitrary retained assumptions.
\end{titlepage}
\tableofcontents
\clearpage

\section{Result and comparison with the current proposal}

The useful simplification is the following single arithmetic inference.
Write \(\overline I\) for the canonical iterator constructor, which
returns its initializer at the literal counter zero and otherwise retains
an iterator node:
\[
 \boxed{\quad
 \frac{\overline I(c,a;u.h)=t\qquad I(Sc,a;u.h)=r}
      {r=N(h[t/u])}
 \quad\text{\rm Step}. \quad}
\]
Here \(N\) contracts zero-counter iterators recursively, and
substitution is capture-free. The rule reads its counter, initializer,
and body from the second premise. Its first premise is required to
have exactly the corresponding predecessor source. No arbitrary
equality template or occurrence path is supplied in the certificate.

In the frozen repository, \texttt{NativeEval.lean} implements the
six-rule computational-equality presentation with separate
\texttt{evalId}, \texttt{evalZero}, and \texttt{evalSucc} constructors
for its three Eval alternatives. The repository also contains
zero-normalization and initializer-only equality proposals.
The incoming \texttt{ipra\_clausal\_canonical.zip} already has a
five-rule clausal calculus, and
\texttt{ipra\_scalar\_uniqueness.zip} already replaces induction by
uniqueness of iteration. Both archives were inspected before selecting
this construction. The present proposal does not claim those earlier
ideas as new.

The new observation is that on zero-normal initialized iterators,
\emph{the successor alternative alone recovers the missing identity
alternative}. For \(J_s=I(S0,s;v.v)\), two Steps give
\[
 \frac{s=t\qquad
       \displaystyle\frac{s=r\qquad J_s=J_s}{J_s=r}\ {\rm Step}}
      {r=t}\ {\rm Step}.
 \tag{1}\label{eq:headline-pivot}
\]
This equality-composition rule works under arbitrary retained assumptions.
It does not depend on induction, numerical cases, or the logical binder.
Together with one-step application of a numerical template, it restores
the entire equality interface.

\subsection{Precise claims}

\begin{theorem}[Main result]\label{thm:main}
Let \(C_{\rm step}\) be the five-rule calculus specified in
Section~\ref{sec:five-pattern-clauses}. It is a faithful presentation of
a theory bi-interpretable with full first-order iPRA.
The interpretation domains are the whole numerical domains, interpreted
equality is actual equality, and both composite interpretations have
isomorphism graph \(x=y\). It is consequently equiconsistent with iPRA.

The new arithmetic inference interface is interderivable, under arbitrary
side assumptions, with normalized successor computation and
one-template equational transport. Replacing the latter pair by Refl
and Step therefore preserves the consequence relation, not merely the
set of true standard numerical equations.
\end{theorem}

The proof has three layers. First, the equality argument
(\ref{eq:headline-pivot}) recovers computation and transport.
Second, the clause interface recovers the native conditional-closure
binder and both case and progress introductions. A nonliteral canonical
term provably equal to zero handles the one delicate introduction case.
Third, the existing internal arithmetic bridge transfers the result to
iPRA. The first two layers are proved in full below; the exact inherited
bridge proofs are reproduced in the appendices with attribution.

\subsection{What the simplicity claim establishes}

There are several distinct comparisons, and they should not be
collapsed into one number:
\begin{itemize}
\item Against the current three-alternative Eval rule, Step has one
  arithmetic pattern. The zero equation is accounted for by the
  existing canonical syntax.
\item Against a matched zero-normal clausal calculus, the number of
  inference schemes remains five. Its explicit computation axiom and
  template-taking transport rule are replaced by Refl and Step.
  The certificate no longer carries either a computation triple or a
  transport template for those steps.
\item Against the incoming five-rule clausal canonical checker, the new
  complete checker is smaller in all three reported source coordinates.
  This comparison also changes the numerical representation, so the
  matched control is the cleaner measurement of the arithmetic change.
\end{itemize}

This report does not beat every historical source-size profile.
The repository contains an aggressively encoded 516-byte guarded
profile, and the incoming uniqueness package has an 866-byte fully
validated implementation of a different certificate protocol. The new
readable checker is not smaller than those programs. There is no
representation-independent theorem here about the globally simplest
axiomatization. The established result is the exact rule simplification
and its controlled complete-checker reduction.

\subsection{Attribution and proof status}

The scalar iterator, the \(D\) binder, clause contexts, the
witness-directed introduction, and zero-normalization are inherited
from the supplied repository. The contribution of this package is the
uniform Step reduction, its combination with the existing logical
interface, the written equivalence proof, and its executable
implementation and translations.

The baseline's Lean modules establish the inherited native/scalar/iPRA
bridge. We inspected those sources but did not compile a new Lean
formalization in this package. Passing Python certificates supplies
implementation evidence; the universal metatheorem is the written
proof, and the two kinds of evidence are kept separate.

